\begin{dfn}\label{def:separation-axioms}
    Let \( X \) be a topological space:
    \begin{itemize}
        \item \( X \) is called a \vocab{\( T_{0} \) space } or a \vocab{Kolmogorov space} if for any pair of points \( x,y \in X \) there exists an open set \( U \) that contains one point but not the other.
        \item \( X \) is called a \vocab{\( T_{1} \) space} if for every pair of points \( x,y \in X \), there exists open sets \( U \) and \( V \) such that \( x \in U  \) but \( x \not \in V \) and \( y \in V \) but \( y \not \in U \).
        \item \( X \) is called a \vocab{\( T_{2} \) space} or a \vocab{Hausdorff space} if for every pair of points \( x,y \in X \), there exists open neighborhood \( U \) containing \( x \) and \( V \) containing \( y \) for which \( U \cap V = \varnothing \). 
        \item \( X \) is called a \vocab{\( T_{3} \) space} or a \vocab{regular space}  if for every point \( x \) and closed set \( C \), there exists a neighborhood \( U \) of \( x \) and an open set \( V \) containing \( C \), for which \( U \cap V =\varnothing \).
        \item \( X \) is called a \vocab{\( T_{4} \) space} or a \vocab{normal space} if for every pair of disjoint closed sets \( C \) and \( D \), there exists open sets \( U \), containing \( C \) and \( V \) containing \( D \), for which \( U \cap V = \varnothing \).
    \end{itemize}
\end{dfn}